\begin{abstract}
\thispagestyle{empty}

Наименование работы: «Веб и мобильное приложение для корпоративного 
обучения с поддержкой адаптивного тестирования».

Цель работы: разработка клиент-серверного программного комплекса 
корпоративного обучения с адаптивным тестированием и логической 
изоляцией данных организаций-арендаторов.

Задачи: анализ предметной области и аналогов; формирование требований; 
проектирование архитектуры и базы данных; реализация серверной части, 
веб-панели и мобильного приложения; описание алгоритмов адаптивного 
тестирования; подготовка программной документации.

Объект и предмет исследования: процессы организации корпоративного 
обучения и методы проектирования клиент-серверных информационных систем 
с адаптивным тестированием.

Работа содержит введение, три главы, заключение, 55 источников 
(25 печатных, 30 электронных) и пять приложений. Объём основного текста — 64 страницы, 
включая 19 рисунков и 14 таблиц.

Во Введении обоснована актуальность, сформулированы цель и задачи, 
определены научная новизна и практическая значимость.

В Первой главе — анализ предметной области, обзор и сравнение 
LMS-аналогов, выбор технологического стека.

Во Второй главе — архитектура системы, модель данных, алгоритмы 
адаптивного тестирования, реализация серверной и клиентской частей.

В Третьей главе — внедрение, тестирование, информационная безопасность 
и надёжность.

В Заключении — выводы и направления развития.

Работоспособность программного комплекса подтверждена опытным развёртыванием на арендованном виртуальном выделенном сервере и публикацией исходного кода в открытом репозитории GitHub.

\end{abstract}
